\documentclass[12pt,a4paper]{report}

\usepackage{polyglossia}
\setmainlanguage{french}

\usepackage{amsmath, amssymb}

\usepackage{graphicx}
\usepackage{caption}

\usepackage{siunitx} % Pour les unités

\usepackage{geometry}
\geometry{top=2cm, bottom=2cm, left=2cm, right=1.5cm}
\setlength{\parindent}{0pt}

\usepackage{enumitem}

\usepackage{hyperref}

\usepackage{xcolor}

\usepackage{multicol}
\setlength{\columnseprule}{1pt}
\def\columnseprulecolor{\color{gray}}

\usepackage{tcolorbox}


\usepackage{import}
\usepackage{xifthen}
\usepackage{pdfpages}
\usepackage{transparent}
\usepackage{float}






\usepackage{myunit}
\usepackage{stanli}














\newcommand{\incfig}[1]{%
	\def\svgwidth{\columnwidth}
	\import{./figures/}{#1.pdf_tex}
}








\newcommand{\RomanNum}[1] %pour les chiffre romain
{\MakeUppercase{\romannumeral #1}}


%commandes pour sigma
\newcommand{\sgma}[1]{\sigma_{#1}}
\newcommand{\sigmax}{\sigma_{\max}}
\newcommand{\sigmin}{\sigma_{\min}}

\newcommand{\sigmabar}[1]{\overline{\sigma}_{#1}}

%fonction pour tau
\newcommand{\taubar}[1]{\overline{\tau}_{#1}}





%fonction unique pour ce document
%force horizontale
\newcommand{\fH}{F_{H}}
\newcommand{\fht}{F_{ht_{1}}}
\newcommand{\fhtt}{F_{ht_{2}}}
\newcommand{\fhq}{F_{hq}}
%moment des fh
\newcommand{\mht}{M_{ht_{1}}}
\newcommand{\mhtt}{M_{ht_{2}}}
\newcommand{\mhq}{M_{hq}}


%force verticale
\newcommand{\fV}{F_{V}}
\newcommand{\fvr}{F_{vr}}
\newcommand{\fvs}{F_{vs_{1}}}
\newcommand{\fvss}{F_{vs_{2}}}
\newcommand{\fvt}{F_{vt_{1}}}
\newcommand{\fvtt}{F_{vt_{2}}}
\newcommand{\fvq}{F_{vq}}
%moment des fv
\newcommand{\mvr}{M_{vr}}
\newcommand{\mvs}{M_{vs_{1}}}
\newcommand{\mvss}{M_{vs_{2}}}
\newcommand{\mvt}{M_{vt_{1}}}
\newcommand{\mvtt}{M_{vt_{2}}}
\newcommand{\mvq}{M_{vq}}


%gamma
\newcommand{\gammah}{\gamma_{h}}
\newcommand{\gammab}{\gamma_{ba}}















\title{Mur de Soutènement}
\author{}
\date{\today}

\begin{document}
	\counterwithin{section}{chapter} % pour que \section{} recommence de 1 aprés \part{}
	
	
	\tableofcontents
	
	\newpage
	
	
	
	
	
	
	
	
	
	\chapter{Vérification de stabilité pour le mur 1}
	
	\begin{tcolorbox}[colback=gray!20, colframe=black, title=Remarque]
		\begin{itemize}
			\item Le calcul se fera pour une bande de 1m.
			\item Le rideau sera calculer comme une console encastrée dans la semelle.
			\item Les semelles avant et arrière seront considérés comme des consoles encastrées dans le rideau.
		\end{itemize}
	\end{tcolorbox}
	
	
	\begin{tcolorbox}[colback=gray!20, colframe=black, title=Remarque]
		La poussée des terres est une action permanente, on doit donc lui affecter une coefficient de 1,35
		
		Tandis que la surcharge est variable donc on lui affecte 1,5.
		
		Toutefois ces coefficient ne sont pas à prendre en considération lorsqu'ils agissent dans un sens favorables.
	\end{tcolorbox}
	
	
	\section{Calcul des forces horizontales}
	
	\subsection{Force de poussée des terres}
	
	\subsubsection{Coefficient K}
	
	\begin{minipage}{.8\textwidth}
		\[ K_{a} = \tan^{2}\left( \dfrac{\pi}{4} - \dfrac{\varphi_{1}}{2} \right)\rightarrow\text{Partie arrière(poussée)} \]
		
		\[ K_{p} = \tan^{2}\left( \dfrac{\pi}{4} + \dfrac{\varphi_{2}}{2} \right)\rightarrow\text{Partie avant(butée)} \]
	\end{minipage}
	\begin{minipage}{.2\textwidth}
		\begin{figure}[H]
			\begin{tikzpicture}
				\draw (0,0) -- (2,0) -- (2,.5) -- (1.3,.5) -- (1.3,2) -- (.8,2) -- (.5,.5) -- (0,.5) -- cycle;
				\node[left] at (.5,1) {Avant};
				\node[right] at (1.3,1) {Arrière};
			\end{tikzpicture}
		\end{figure}
	\end{minipage}
	
	\begin{tcolorbox}[colback=red!20, colframe=red, title=Remarque]
		ATTENTION AU UNITÉ DE $\varphi$ ET DE $\pi$.
		
		C'est à dire si $\varphi$ est en degré alors $\pi$ doit aussi être en degré.
		
		Généralement les deux sont en degré.
	\end{tcolorbox}
	
	
	\[ K = K_{a} = K_{b} = 0,271 \]
	
	\vspace{\baselineskip}
	
	
	
	\begin{tcolorbox}[colback=white, colframe=gray!50, title=\textcolor{black}{\textbf{Les forces}}]
		
		\begin{figure}[H]
			\centering
			\begin{tikzpicture}
				\point{a}{0}{0};
				\point{f}{0}{.6};
				\point{b}{0}{2};
				
				\lineload{1}{a}{b}[2][0];
				
				\draw[green,<-,very thick] (f) -- ++(180:3)node[left] {$F_{ht}$};
			\end{tikzpicture}
		\end{figure}
		
		\begin{multicols}{2}
			Terres situé à l'avant
			
			\[  
			\begin{aligned}
				F_{ht_{2}} &= K_{a} \times \dfrac{(\gamma_{h} \times h_{2}) \times h_{2}}{2} \\
				&= 0,271 \times \dfrac{19,8 \times 1,3^{2}}{2} \\
				&= 4,534\kn 
			\end{aligned}
			\]
			
			\columnbreak
			Terre situé à l'arrière
			\[
			\begin{aligned}
				\fht &= K \times \dfrac{(\gammah \times H) \times H}{2} \\
				&= 0,271 \times \dfrac{19,8 \times 4,8^{2}}{2} \\
				&= 61,814\kn
			\end{aligned}
			\]
		\end{multicols}
	\end{tcolorbox}
	
	
	
	
	
	\begin{tcolorbox}[colback=white, colframe=gray!50, title=\textcolor{black}{\textbf{Bras de levier}}]
		
		\begin{figure}[H]
			\centering
			\begin{tikzpicture}
				\draw (0,0) node[red,below] {A} -- (2,0) -- (2,.5) -- (1.3,.5) -- (1.3,2) -- (.8,2) -- (.5,.5) -- (0,.5) -- cycle;
				\draw[red,fill] (0,0) circle (0.1);
			\end{tikzpicture}
		\end{figure}
		
		\begin{multicols}{2}
			Avant :
			\[ B_{L} = h_{2}/3 = 1,3/3 = 0,433\meterl \]
			
			\columnbreak
			Arrière :
			
			\[ B_{L} = \dfrac{h_{1}}{3} = \dfrac{4,8}{3} = 1,6\meterl \]
		\end{multicols}
	\end{tcolorbox}
	
	
	\begin{tcolorbox}[colback=white, colframe=gray!50, title=\textcolor{black}{\textbf{Calcul du moment par rapport au point "A"}}]
		\begin{multicols}{2}
			Avant :
			\[
			\begin{aligned}
				\mhtt &= \fhtt \times \dfrac{h_{2}}{3} \\
				&= 4,534 \times 0,433 \\
				&= 1,964\knxm
			\end{aligned}
			\]
			
			\columnbreak
			Arrière : 
			\[
			\begin{aligned}
				\mht &= \fht \times B_{L}\\
				&= 61,814 \times 1,6 \\
				&= 98,902\knxm 
			\end{aligned}
			\]
		\end{multicols}
	\end{tcolorbox}
	
	
	\subsection{Force de poussée due à la surcharge}
	
	
	\begin{tcolorbox}[colback=white, colframe=gray!50, title=\textcolor{black}{\textbf{La force}}]
		\[ \fhq = K \times q \times H = 0,271 \times (4 \times 4,8) = 4,752\kn \]
	\end{tcolorbox}
	
	\begin{tcolorbox}[colback=white, colframe=gray!50, title=\textcolor{black}{\textbf{Bras de levier}}]
		\[ B_{L} = \dfrac{H}{2} = \dfrac{4,8}{2} = 2,4\meterl \]e
	\end{tcolorbox}
	
	\begin{tcolorbox}[colback=white, colframe=gray!50, title=\textcolor{black}{\textbf{Le moment/A}}]
		\[ \mhq = \fhq \times B_{L} = 4,752 \times 2,4 = 11,405\knxm \]
	\end{tcolorbox}
	
	
	\section{Calcul des forces verticales}
	
	\subsection{Poids propre du rideau}
	
	
	\begin{minipage}{.5\textwidth}
		\begin{figure}[H]
			\centering
			\begin{tikzpicture}
				\point{A}{-3}{0};
				\point{a}{0}{0};
				\point{b}{2}{0};
				\point{c}{2}{2.5};
				\point{d}{1}{2.5};
				\point{g}{1}{1};
				
				\dimensioning{1}{a}{b}{-.7}[$E_{b}$];
				\dimensioning{1}{a}{g}{-1.5}[$2/3E_{b}$];
				\dimensioning{1}{A}{a}{-1.5}[$b_{2}$];
				\dimensioning{1}{d}{c}{3}[$E_{h}$];
				
				\dimensioning{2}{a}{g}{-.5}[$V_{1}$];
				\dimensioning{2}{g}{d}{-.5}[$V_{2}$];
				\dimensioning{2}{c}{b}{2.5}[H];
				
				\draw (a) -- (b) -- (c) -- (d) -- cycle;
				\node[circle,draw] at (g) {\tiny G};
				\node[red,circle,draw] at (A) {A};
			\end{tikzpicture}
		\end{figure}
	\end{minipage}
	\begin{minipage}{.5\textwidth}
		\[ S = \dfrac{(E_{h} + E_{b})}{2} \times H = 0,35\msq \]
		
		\[ V_{1} = \dfrac{H}{3} \times \dfrac{(2E_{b} + E_{h})}{E_{b} + E_{h}} \]
		\[ V_{2} = \dfrac{H}{3} \times \dfrac{(E_{b} + 2E_{h})}{E_{b} + E_{h}} \]
		
	\end{minipage}
	
	
	\begin{tcolorbox}[colback=white, colframe=gray!50, title=\textcolor{black}{\textbf{La force}}]
		\[ \begin{aligned}
			\fvr &= \gammah \times S \times H \\
			&= 25 \times 0,35 \times 4,8 \\
			&= 42\kn
		\end{aligned} \]
	\end{tcolorbox}
	
	\begin{tcolorbox}[colback=white, colframe=gray!50, title=\textcolor{black}{\textbf{Bras de levier}}]
		\[ B_{L} = b_{2} + \dfrac{2}{3} \times E_{b} = 1 + \dfrac{2}{3} \times 0,4 = 1,266\meterl \]
	\end{tcolorbox}
	
	\begin{tcolorbox}[colback=white, colframe=gray!50, title=\textcolor{black}{\textbf{Le moment/A}}]
		\[ \mvr = \fvr \times B_{L} = 42 \times 1,266 = 53,17\knxm \]
	\end{tcolorbox}
	
	\subsection{Poids propre des semelles}
	
	\begin{tcolorbox}[colback=white, colframe=gray!50, title=\textcolor{black}{\textbf{Schémas}}]
		\begin{multicols}{2}
			\centering
			Semelle avant :
			
			\begin{tikzpicture}
				\draw (0,0) node[red,below] {A} coordinate (a) rectangle (2.5,.5) coordinate (b);
				\draw[very thick,green,->] (1.25,.25) -- ++(-90:1) node[below]{$\fvss$};
				\dimensioning{1}{a}{b}{1}[$b_{2}$];
				\dimensioning{2}{a}{b}{3}[$E_{s}$];
			\end{tikzpicture}
			
			\columnbreak
			Semelle arrière : 
			
			\begin{tikzpicture}
				\draw (0,0) coordinate (a) rectangle (2.5,.5) coordinate (b);
				\draw[very thick,green,->] (1.25,.25) -- ++(-90:1) node[below]{$\fvss$};
				\draw[dashed,black!40] (-2,0) node[red,below] {A} rectangle (0,.5);
				\coordinate (A) at (-2,0);
				\coordinate (oa) at (-1,0);
				\dimensioning{1}{a}{b}{1}[$b_{1}$];
				\dimensioning{2}{a}{b}{3}[$E_{s}$];
				\dimensioning{1}{A}{oa}{1}[$b_{2}$];
				\dimensioning{1}{oa}{a}{1}[$E_{b}$];
			\end{tikzpicture}
		\end{multicols}
	\end{tcolorbox}
	
	
	\begin{tcolorbox}[colback=white, colframe=gray!50, title=\textcolor{black}{\textbf{La force}}]
		\begin{multicols}{2}
			Avant : 
			
			\[
			\begin{aligned}
				\fvss &= \gammab \times b_{2} \times E_{s} \\
				&= 25 \times 1 \times 0,4 \\
				&= 10\kn 
			\end{aligned}
			\]
			
			\columnbreak
			
			\[ \begin{aligned}
				\fvs &= \gammab \times b_{1} \times E_{s} \\
				&= 25 \times 1,5 \times 0,4 \\
				&= 15\kn
			\end{aligned} \]
			
			
		\end{multicols}
	\end{tcolorbox}
	
	\begin{tcolorbox}[colback=white, colframe=gray!50, title=\textcolor{black}{\textbf{B.L}}]
		\begin{multicols}{2}
			\[ B_{L} = \dfrac{b_{2}}{2} = \dfrac{1}{2} = 0,5\meterl \]
			\columnbreak
			\[ \begin{aligned}
				B_{L} &= b_{2} + E_{b} + \dfrac{b_{1}}{2} \\
				&=1 + 0,4 + \dfrac{1,5}{2} \\
				&=2,15\meterl
			\end{aligned}\]
		\end{multicols}
	\end{tcolorbox}
	
	\begin{tcolorbox}[colback=white, colframe=gray!50, title=\textcolor{black}{\textbf{Moment/A}}]
		\begin{multicols}{2}
			\[ \begin{aligned}
				\mvss &= \fvss \times B_{L} \\
				&= 10 \times 0,5 \\
				&= 5\knxm
			\end{aligned} \]
			
			\columnbreak
			
			\[ \begin{aligned}
				\mvs &= \fvs \times B_{L} \\
				&= 15 \times 2,15 \\
				&= 32,25\knxm
			\end{aligned} \]
			
		\end{multicols}
	\end{tcolorbox}
	
	
	\subsection{Poids des terres}
	
	
	
	\begin{tcolorbox}[colback=white, colframe=gray!50, title=\textcolor{black}{\textbf{Schéma}}]
		\centering
		\begin{tikzpicture}[scale=2]
			\draw (0,0) -- (2,0) -- (2,.5) -- (1.3,.5) -- (1.3,2) -- (.8,2) -- (.5,.5) -- (0,.5) -- cycle;
			\draw[pattern=north west lines] (0,.5) rectangle (.8,1);
			\draw[pattern=north west lines] (1.3,.5) rectangle (2,2);
			
			\coordinate (ft) at ($(0,.5)!.5!(.8,1)$);
			\coordinate (f) at ($(1.3,.5)!.5!(2,2)$);
			\draw[very thick,green,->] (ft) -- ++(-90:1.3) node[below]{$\fvtt$};
			\draw[very thick,green,->] (f) -- ++(-90:2) node[below]{$\fvt$};
		\end{tikzpicture}
	\end{tcolorbox}
	
	
	
	\begin{tcolorbox}[colback=white, colframe=gray!50, title=\textcolor{black}{\textbf{La force}}]
		\begin{multicols}{2}
			Avant :
			\[ \begin{aligned}
				\fvtt &= \gammah \times b_{2} \times h_{2}' \\
				&= \gammah \times b_{2} \times (h_{2} - E_{s}) \\
				&= 19,8 \times 1 \times (1,3 - 0,4) \\
				&= 17,82\kn
			\end{aligned} \]
			
			\columnbreak
			
			Arrière : 
			\[ \begin{aligned}
				\fvt &= \gammah \times b_{1} \times h_{1}' \\
				&= \gammah \times b_{1} \times (H - E_{s}) \\
				&= 19,8 \times 1,5 \times (4,8 - 0,4) \\
				&= 130,68\kn
			\end{aligned} \]
			
		\end{multicols}
	\end{tcolorbox}
	
	
	\begin{tcolorbox}[colback=white, colframe=gray!50, title=\textcolor{black}{\textbf{B.L}}]
		Avant :
		\begin{multicols}{2}
			\[ \begin{aligned}
				B_{L} = \dfrac{b_{2}}{2} = 1/2 = 0,5\meterl
			\end{aligned} \]
			
			\columnbreak
			
			Arrière :
			\[ \begin{aligned}
				B_{L} = b_{2}  + E_{b} + \dfrac{b_{1}}{2} = 2,15\meterl
			\end{aligned} \]
		\end{multicols}
	\end{tcolorbox}
	
	\begin{tcolorbox}[colback=white, colframe=gray!50, title=\textcolor{black}{\textbf{Moment/A}}]
		\begin{multicols}{2}
			Avant :
			\[ \begin{aligned}
				\mvtt &= \fvtt \times B_{L} \\
				&= 17,82 \times 0,5 \\
				&= 8,91\knxm
			\end{aligned} \]
			
			\columnbreak
			
			Arrière :
			\[ \begin{aligned}
				\mvt &= \fvt \times B_{L} \\
				&= 130,68 \times 2,15 \\
				&= 280,962\knxm
			\end{aligned} \]
		\end{multicols}
	\end{tcolorbox}
	
	
	
	\subsection{Surcharge}
	
	\begin{tcolorbox}[colback=white, colframe=gray!50, title=\textcolor{black}{\textbf{La force}}]
		\begin{figure}[H]
			\centering
			\begin{tikzpicture}
				\draw (0,0) -- (2,0) -- (2,.5) -- (1.3,.5) -- (1.3,2) coordinate (e) -- (.8,2) coordinate (f) -- (.5,.5) -- (0,.5) -- cycle;
				\draw[pattern=north east lines] (e) rectangle (2,2.5) node[midway,above=.2] {q};
				\coordinate (f) at ($(1.3,2)!.5!(2,2)$);
				\draw[very thick,green,->] (f) -- ++(-90:1) node[right]{$\fvq$};
			\end{tikzpicture}
		\end{figure}
		
		\[ \fvq = q \times b_{1} = 4 \times 1,5 = 6\kn \]
	\end{tcolorbox}
	
	\begin{tcolorbox}[colback=white, colframe=gray!50, title=\textcolor{black}{\textbf{B.L}}]
		\[ B_{L} = b_{2} + E_{b} + \dfrac{b_{1}}{2} = 2,15\meterl \]
	\end{tcolorbox}
	
	\begin{tcolorbox}[colback=white, colframe=gray!50, title=\textcolor{black}{\textbf{Moment/A}}]
		\[ \mvq = \fvq \times B_{L} = 6 \times 2,15 = 12,9\knxm \]
	\end{tcolorbox}
	
	
	
	\section{Stabilité au glissement}
	
	\[ \dfrac{\sum \fH}{\sum \fV} \leq 0,6 \]
	
	\[ \begin{aligned}
		\sum \fH &= 1,35 \times (\fht - \fhtt) + 1,5 \times \fhq \quad \text{(défavorable)} \\
		&= 1,35 \times (61,814 - 4,534) + 1,5 \times 4,752 \\
		&= 84,456\kn 
	\end{aligned} \]
	
	\[ \begin{aligned}
		\sum \fV &= \fvr + \fvs + \fvss + \fvt + \fvtt + \fvq \quad \text{(favorable)} \\
		&= 42 + 10 + 15 + 17,82 + 130,68 + 6 \\
		&= 221,5\kn
	\end{aligned} \]
	
	
	\[ \dfrac{\sum \fH}{\sum \fV} = \dfrac{84,456}{221,5} = 0,381 < 0,6 \]
	
	{\centering Pas de risque de glissement\par}
	
	\begin{tcolorbox}[colback=gray!20, colframe=black, title=Remarque]
		Si cette condition ne passe pas soit on augmente les dimensions du mur ou on ajoute une bêche.
	\end{tcolorbox}
	
	
	
	\section{Stabilité au renversement}
	
	\[ \dfrac{\sum M_{stabilisant}}{\sum M_{renversant}} \geq 1,5 \]
	
	\[ \begin{aligned}
		\sum M_{renversant} &= 1,35 \times (\mht - \mhtt) + 1,5 \times \mhq \\
		&= 147,975\knxm
	\end{aligned} \]
	
	\[ \begin{aligned}
		\sum M_{stabilisant} &= \mvr + \mvt + \mvtt + \mvs + \mvss + \mvq \\
		&= 393,19\knxm 
	\end{aligned} \]
	
	\[ \dfrac{\sum M_{stabilisant}}{\sum M_{renversant}} = \dfrac{393,194}{149,975} = 2,657 > 1,5 \]
	
	{\centering Le mur est stable vis-à-vis au renversement\par}
	
	\section{Stabilité au poinçonnement}
	
	Il faut ramener tous les efforts au CDG de \textcolor{red}{la semelle}
	
	\subsection{Vérification à ELS}
	
	\[ M_{s} = M_{s1}^{retient} - M_{s2}^{renverse} \]
	
	
	\[ M_{s1} = \sum M_{stabilisant} = 393,19\knxm \quad \text{(ELS)} \]
	
	\[ M_{s2} = \sum M_{renversant} = 108,344\knxm \quad \text{(ELS)} \]
	
	
	\[ M_{s} = 393,194 - 108,344 = 284,85\knxm \]
	
	\vspace{\baselineskip}
	
	\[ M_{g} = \sum \fV \times e \]
	
	\[ e = \dfrac{B}{2} - \dfrac{M_{s}}{\sum \fV} = \dfrac{2,9}{2} - \dfrac{284,85}{221,5} = 0,164\meterl \]
	
	\[ M_{g} = 221,5 \times 0,164 = 36,325\knxm \]
	
	
	\[ \begin{cases}
		\sgma{\max} = \dfrac{N}{S} + \dfrac{M \times V}{I} = \dfrac{\sum \fV}{B \times 1ml} + \dfrac{M_{g} \times \dfrac{B}{2}}{\dfrac{1 \times b^{3}}{12}} = 0,102\mpa \\ \\
		\sgma{\min} = \dfrac{N}{S} - \dfrac{M \times V}{I} = \dfrac{\sum \fV}{B \times 1ml} - \dfrac{M_{g} \times \dfrac{B}{2}}{\dfrac{1 \times b^{3}}{12}} = 0,05\mpa
	\end{cases} \]
	
	$B = 2,9\meterl$
	
	
	
	\begin{figure}[H]
		\centering
		\begin{tikzpicture}[scale=2]
			\draw (0,0) -- (2,0) -- (2,.5) -- (1.3,.5) -- (1.3,2) -- (.8,2) -- (.5,.5) -- (0,.5) -- cycle;
			\draw[pattern=vertical lines] (0,-.2) -- (2,-.2) -- (2,-.7) node[below] {$\sgma{\min}$} -- (0,-1.2) node[below] {$\sgma{\max}$} -- cycle;
		\end{tikzpicture}
	\end{figure}
	
	
	\[ \sgma{moy} = \dfrac{3\sgma{\max} + \sgma{\min}}{4} = 0,089\mpa \]
	
	\[ \sgma{moy} = 0,089 < \sgma{sol} = 0,2\mpa \]
	
	{\centering Condition vérifiée \par}
	
	\vspace{\baselineskip}
	On refait le calcul a l'ELU mais on applique les coefficient uniquement pour les forces défavorable.
	
	
	
	
	
	
	
	
	
	
	
	
	
	
	
	
	
	\vfill \hfill \Huge \LaTeX
	
\end{document}